\chapter{dg-modules}
\label{ch:modules}

This chapter is Keller \S 3.1. Everything here is stated for right dg-modules; left
dg-modules are right dg-modules over $\cA^{\op}$.

\begin{definition}[Right dg-modules]
\label{def:dg-module}
% planned: DgCat.DGModule
\uses{def:dg-functor, def:dg-op, ex:dg-complexes}
A right dg-module over a small dg-category $\cA$ is a dg-functor
$M : \cA^{\op} \to \Chdg(k)$. Equivalently, it is a family of complexes $M(X)$ together with
chain maps $M(Y) \otimes \cA(X,Y) \to M(X)$ compatible with composition and units.
\end{definition}

\begin{definition}[The dg-category of dg-modules]
\label{def:dg-module-dgcat}
% planned: DgCat.DGModule.dgCategory
\uses{def:dg-module, def:dg-functor-category}
$\Chdg(\cA) = \Hom(\cA^{\op}, \Chdg(k))$, the dg-category of dg-functors of
Definition~\ref{def:dg-functor-category}. We write $\Hom(L,M)$ for its hom-complexes.
\end{definition}

\begin{definition}[The category of dg-modules]
\label{def:dg-module-cat}
% planned: DgCat.DGModule.category
\uses{def:dg-module-dgcat, def:Z0-H0}
$\Ch(\cA) = \Zz(\Chdg(\cA))$: objects are the dg-modules, morphisms are the morphisms of
dg-functors.
\end{definition}

\begin{proposition}[$\Ch(\cA)$ is abelian]
\label{prop:dg-module-abelian}
% planned: DgCat.DGModule.abelian
\uses{def:dg-module-cat}
$\Ch(\cA)$ is an abelian category; a morphism $L \to M$ is a monomorphism (epimorphism) iff
each component $L(X)^n \to M(X)^n$ is injective (surjective). Kernels and cokernels are
computed objectwise and degreewise.
\end{proposition}
\begin{proof}
Limits and colimits of dg-modules are computed in $\Ch(k)$ objectwise; the dg-module
structure maps are inherited. Compare mathlib's proof that functor categories into an
abelian category are abelian.
\end{proof}

\begin{definition}[Homology of a dg-module]
\label{def:dg-module-homology}
% planned: DgCat.DGModule.homologyFunctor
\uses{def:dg-module}
$H^n(M)$ is the functor $X \mapsto H^n(M(X))$; a morphism of dg-modules is a
quasi-isomorphism if it induces isomorphisms on all $H^n(M(X))$.
\end{definition}

\begin{definition}[Representable dg-modules]
\label{def:dg-representable}
% planned: DgCat.DGModule.representable
\uses{def:dg-module}
For an object $X$ of $\cA$, $X^{\wedge} = \cA(?,X)$ is the right dg-module represented by $X$.
\end{definition}

\begin{lemma}[dg-Yoneda]
\label{lem:dg-yoneda}
% planned: DgCat.DGModule.yoneda
\uses{def:dg-representable, def:dg-module-dgcat}
For a dg-module $M$ and an object $X$ there is a natural isomorphism of complexes
$\Hom(X^{\wedge}, M) \xrightarrow{\sim} M(X)$ (Keller (1)).
\end{lemma}
\begin{proof}
The enriched Yoneda lemma, evaluated at the unit.
\end{proof}

\begin{definition}[Shift of dg-modules]
\label{def:dg-module-shift}
% planned: DgCat.DGModule.shift
\uses{def:dg-module}
$M[n]$ is the dg-module $Y \mapsto M(Y)[n]$, with the shifted differential
$(-1)^n d$ on each $M(Y)$.
\end{definition}

\begin{definition}[dg-modules up to homotopy]
\label{def:dg-module-homotopy-cat}
% planned: DgCat.DGModule.homotopyCategory
\uses{def:dg-module-dgcat, def:Z0-H0}
$\Hcat(\cA) = \Hz(\Chdg(\cA))$, the category of dg-modules up to homotopy, with the
canonical functor $\Ch(\cA) \to \Hcat(\cA)$.
\end{definition}

\begin{lemma}[Homotopy classes out of representables]
\label{lem:dg-homotopy-yoneda}
% planned: DgCat.DGModule.homotopyYoneda
\uses{lem:dg-yoneda, def:dg-module-shift, def:dg-module-homotopy-cat}
For $n \in \mathbb{Z}$ there are natural isomorphisms
$\Hcat(\cA)(X^{\wedge}, M[n]) \cong H^n \Hom(X^{\wedge}, M) \cong H^n M(X)$ (Keller (2)).
\end{lemma}
\begin{proof}
\uses{lem:dg-yoneda}
Take $H^0$ of the isomorphism of Lemma~\ref{lem:dg-yoneda} after shifting.
\end{proof}

\begin{proposition}[dg-modules over an algebra]
\label{prop:dg-module-ring}
% planned: DgCat.DGModule.ringComparisonC, DgCat.DGModule.ringComparisonH
\uses{ex:dg-algebra, def:dg-module-cat, def:dg-module-homotopy-cat, lem:Z0-H0-complexes}
If $\cA = \cA_B$ is the one-object dg-category of a $k$-algebra $B$, then
$\Ch(\cA) \cong \Ch(B)$, $\Chdg(\cA) \cong \Chdg(B)$ and $\Hcat(\cA) \cong \Hcat(B)$,
and $X^{\wedge}$ for the unique object $X$ is the free module $B$ in degree $0$.
\end{proposition}
\begin{proof}
\uses{lem:Z0-H0-complexes}
A dg-functor out of a one-object dg-category concentrated in degree $0$ is a complex with a
compatible $B$-action; unwind and apply Lemma~\ref{lem:Z0-H0-complexes}.
\end{proof}
